\section*{Capítulo \ref{ch:g12:cells-and-electrolysis} --- Pilhas e eletrólise}

\begin{solution}{exo:g12:cells-and-electrolysis:1}
O eletrodo de zinco é o \omterm{def:g12:cells-and-electrolysis:anode}{ânodo}, \ce{Zn -> Zn^{2+} + 2e-}; o eletrodo de
cobre é o \omterm{def:g12:cells-and-electrolysis:anode}{cátodo}, \ce{Cu^{2+} + 2e- -> Cu}.
\end{solution}

\begin{solution}{exo:g12:cells-and-electrolysis:2}
% ledger: const:F
$Q = 0.20 \times 1800 = \qty{360}{C}$; $n(e^-) = 360/96485 =
\qty{3.7e-3}{mol}$.
\end{solution}

\begin{solution}{exo:g12:cells-and-electrolysis:3}
Ela fecha o circuito e mantém cada solução eletricamente neutra, com os
seus \omterm{def:g9:ions:ion}{íons} entrando nos compartimentos. Sem ela, o circuito fica aberto:
não passa corrente, o voltímetro indica zero e uma lâmpada continua
apagada.
\end{solution}

\begin{solution}{exo:g12:cells-and-electrolysis:4}
$2.5 \times 3600 = \qty{9000}{C}$.
\end{solution}

\begin{solution}{exo:g12:cells-and-electrolysis:5}
Não: a reação vai no sentido oposto ao espontâneo. O gerador fornece a
energia.
\end{solution}

\begin{solution}{exo:g12:cells-and-electrolysis:6}
\omterm{def:g12:cells-and-electrolysis:anode}{Ânodo} (cobre, $-$): \ce{Cu -> Cu^{2+} + 2e-}. \omterm{def:g12:cells-and-electrolysis:anode}{Cátodo} (prata, $+$):
\ce{Ag+ + e- -> Ag}. Global: \ce{Cu + 2Ag+ -> Cu^{2+} + 2Ag}.
\end{solution}

\begin{solution}{exo:g12:cells-and-electrolysis:7}
Três anos são $3 \times 365 \times 24 = \qty{26280}{h}$;
$0.010 \times 26280 = \qty{263}{mA.h}$.
\end{solution}

\begin{solution}{exo:g12:cells-and-electrolysis:8}
% ledger: const:F, aw:Ag
$Q = 0.50 \times 2400 = \qty{1200}{C}$; $n(e^-) = n(\ce{Ag}) =
1200/96485 = \qty{0.0124}{mol}$; $m = 0.0124 \times 107.9 =
\qty{1.34}{g}$.
\end{solution}

\begin{solution}{exo:g12:cells-and-electrolysis:9}
A corrente vai no fio do cobre ($+$) para o zinco ($-$), no sentido
oposto ao dos \omterm{def:g9:inside-the-atom:nucleus}{elétrons}. Os \omterm{def:g9:ions:ion}{íons} \ce{K+} se deslocam para o compartimento
do cobre, que perde \ce{Cu^{2+}}; os \omterm{def:g9:ions:ion}{íons} \ce{NO3-}, para o
compartimento do zinco, que ganha \ce{Zn^{2+}}.
\end{solution}

\begin{solution}{exo:g12:cells-and-electrolysis:10}
Nos dois casos, o \omterm{def:g12:cells-and-electrolysis:anode}{ânodo} é onde acontece a \omterm{def:g11:redox:oxidation}{oxidação}. Numa pilha, o \omterm{def:g12:cells-and-electrolysis:anode}{ânodo}
é $-$, a reação é espontânea e a energia química vira energia elétrica;
numa \omterm{def:g12:cells-and-electrolysis:electrolysis}{eletrólise}, o \omterm{def:g12:cells-and-electrolysis:anode}{ânodo} está ligado ao $+$ do gerador, a reação é
forçada e a energia elétrica vira energia química.
\end{solution}

\begin{solution}{exo:g12:cells-and-electrolysis:11}
\qty{12}{mL}: a equação \ce{2H2O -> 2H2 + O2} dá duas \omterm{def:g7:atoms-and-molecules:molecule}{moléculas} de gás
hidrogênio para uma de gás oxigênio, e quantidades iguais de gás ocupam
volumes iguais.
\end{solution}

\begin{solution}{exo:g12:cells-and-electrolysis:12}
% ledger: aw:Zn, const:F
$n(\ce{Zn}) = 5.0/65.4 = \qty{0.076}{mol}$;
$n(\ce{Cu^{2+}}) = 0.100 \times 0.50 = \qty{0.050}{mol}$; eles reagem um
para um, logo os \omterm{def:g9:ions:ion}{íons} cobre acabam primeiro. $n(e^-) = 2 \times 0.050
= \qty{0.100}{mol}$, $Q = \qty{9.65e3}{C} = \qty{2.68}{A.h}$. A
\qty{50}{mA}: $2.68/0.050 = \qty{54}{h}$.
\end{solution}

\begin{solution}{exo:g12:cells-and-electrolysis:13}
$Q = 2.0 \times 3600 = \qty{7200}{C}$; $E = 7200 \times 1.5 =
\qty{1.08e4}{J}$, isto é, $1.5 \times 2.0 = \qty{3.0}{W.h}$.
\end{solution}

\begin{solution}{exo:g12:cells-and-electrolysis:14}
% ledger: const:F
$Q = \qty{7200}{C}$, $n(e^-) = \qty{0.0746}{mol}$. Dois \omterm{def:g9:inside-the-atom:nucleus}{elétrons} por
\ce{H2}: $n(\ce{H2}) = \qty{0.0373}{mol}$, $V = \qty{0.90}{L}$. Quatro
por \ce{O2}: $n(\ce{O2}) = \qty{0.0187}{mol}$, $V = \qty{0.45}{L}$.
\end{solution}

\begin{solution}{exo:g12:cells-and-electrolysis:15}
O cobre dissolvido no \omterm{def:g12:cells-and-electrolysis:anode}{ânodo} é igual ao cobre depositado, já que os
mesmos \omterm{def:g9:inside-the-atom:nucleus}{elétrons} passam nos dois: $2.84/3.00 = \qty{94.7}{\percent}$ de
cobre.
\end{solution}

\begin{solution}{pb:g12:cells-and-electrolysis:1}
% ledger: const:F, const:NA, aw:Li, dch:CH4
\textbf{1.} O \omterm{def:g12:cells-and-electrolysis:anode}{ânodo}: o lítio é oxidado ali.

\textbf{2.} O eletrodo de grafite, o \omterm{def:g12:cells-and-electrolysis:anode}{ânodo}.

\textbf{3.} Os \omterm{def:g9:inside-the-atom:nucleus}{elétrons} percorrem o celular do eletrodo de grafite para
o eletrodo de óxido; os \omterm{def:g9:ions:ion}{íons} lítio atravessam o eletrólito no mesmo
sentido, da grafite para o óxido.

\textbf{4.} A sua reação pode ser forçada no sentido inverso por um
carregador.

\textbf{5.} \qty{4.000}{A.h}, isto é, $4.000 \times 3600 =
\qty{14400}{C}$.

\textbf{6.} $14400/96485 = \qty{0.1492}{mol}$.

\textbf{7.} $0.1492 \times \num{6.022e23} = \num{8.99e22}$ \omterm{def:g9:inside-the-atom:nucleus}{elétrons}.

\textbf{8.} $4.000/0.25 = \qty{16}{h}$.

\textbf{9.} $0.20 \times 14400 = \qty{2880}{C}$.

\textbf{10.} $E = 14400 \times 3.7 = \qty{5.3e4}{J}$, cerca de
\qty{53}{kJ}.

\textbf{11.} $53280/3600 = \qty{14.8}{W.h}$.

\textbf{12.} Queimar \qty{1}{g} de metano libera um pouco mais de
energia (\qty{55.7}{kJ}) do que uma bateria de celular carregada
armazena.

\textbf{13.} $1000/14.8 = 67.6$: 67 cargas completas.

\textbf{14.} \ce{Li+ + e- -> Li}: uma \omterm{def:g11:redox:oxidation}{redução}.

\textbf{15.} O carregador força a reação no sentido oposto ao
espontâneo.

\textbf{16.} $14400/2.0 = \qty{7200}{s}$, \qty{2.0}{h}.

\textbf{17.} Um \omterm{def:g9:ions:ion}{íon} lítio por \omterm{def:g9:inside-the-atom:nucleus}{elétron}: \qty{0.1492}{mol}.

\textbf{18.} $0.1492 \times 6.9 = \qty{1.03}{g}$.

\textbf{19.} Cerca de \qty{1.03}{g} de lítio vai e vem entre os
eletrodos numa carga completa.
\end{solution}
